\begin{figure}[H]
\centering
\begin{tikzpicture}[
  % every box gets the SAME fixed width and height, so a box with three lines of
  % text cannot grow upwards into the header band above it
  obox/.style={draw=CVBLUE!70!black, fill=CVBLUE!10, rectangle, rounded corners=3pt,
    text width=2.0cm, minimum height=1.60cm,
    align=center, font=\small, inner sep=4pt},
  rbox/.style={draw=black!55, fill=black!5, rectangle, rounded corners=3pt,
    text width=2.0cm, minimum height=1.60cm,
    align=center, font=\small, inner sep=4pt},
  ohdr/.style={draw=CVBLUE!80!black, fill=CVBLUE, rectangle, rounded corners=4pt,
    minimum width=14.6cm, minimum height=0.52cm,
    align=left, font=\small\bfseries, text=white, inner sep=6pt},
  rhdr/.style={draw=black!70, fill=black!70, rectangle, rounded corners=4pt,
    minimum width=14.6cm, minimum height=0.52cm,
    align=left, font=\small\bfseries, text=white, inner sep=6pt},
  arr/.style={-{Stealth[length=2.7mm,width=2.1mm]}, line width=0.9pt,
    CVBLUE!70!black, shorten >=1pt, shorten <=1pt},
  darr/.style={-{Stealth[length=2.7mm,width=2.1mm]}, line width=0.9pt, dashed,
    dash pattern=on 3pt off 2.5pt, black!50},
]

%% ── Handoff drawn FIRST, so the runtime header band paints over the middle of
%%    it and the line reads as passing behind the band rather than through the
%%    title text. It re-emerges just above the box it feeds. ─────────────────
\draw[darr]
  (6.10,-2.05) -- (6.10,-2.50)
  -- node[above, pos=0.62, font=\scriptsize, text=black!60] {trained model files}
  (0.00,-2.50)
  -- (0.00,-3.60);

%% ── Section headers ──────────────────────────────────────────────────────────
\node[ohdr] at (0,  0.00) {(1)\enspace Offline Preparation --- done once};
\node[rhdr] at (0, -3.00) {(2)\enspace Runtime Analysis --- every message the user checks};

%% ── Offline row ──────────────────────────────────────────────────────────────
\node[obox] (dc)  at (-6.10, -1.25) {Seed articles\\+ generated\\messages};
\node[obox] (qrv) at (-3.05, -1.25) {Schema +\\quality\\review};
\node[obox] (spl) at ( 0.00, -1.25) {Split by\\source\\article};
\node[obox] (qlo) at ( 3.05, -1.25) {QLoRA +\\PhoBERT\\training};
\node[obox] (gex) at ( 6.10, -1.25) {Model\\packaging};

\draw[arr] (dc)--(qrv);
\draw[arr] (qrv)--(spl);
\draw[arr] (spl)--(qlo);
\draw[arr] (qlo)--(gex);

%% ── Runtime row ──────────────────────────────────────────────────────────────
\node[rbox] (inp) at (-6.10, -4.40) {\textbf{Input}\\one pasted\\message};
\node[rbox] (rif) at (-3.05, -4.40) {Runtime\\interface};
\node[rbox] (mod) at ( 0.00, -4.40) {Local model\\backend};
\node[rbox] (dly) at ( 3.05, -4.40) {Decision\\layer};
\node[rbox] (out) at ( 6.10, -4.40) {\textbf{Output}\\risk, label,\\cues, advice};

\draw[arr] (inp)--(rif);
\draw[arr] (rif)--(mod);
\draw[arr] (mod)--(dly);
\draw[arr] (dly)--(out);

\end{tikzpicture}
\caption{How the system fits together. The top row is the offline work, done once: crawling the source articles, generating and checking messages, splitting them by source article, training both models, and packaging the result. The dashed line is the only thing that crosses between the two rows --- the trained model files. The bottom row is what happens each time a user checks a message: one pasted message goes in on the left, and a JSON result with the risk tier, the label, the quoted cues, and the advice comes out on the right.}
\label{fig:system-overview}
\end{figure}
